\section{Related work}

%In \cite{moryossef2021evaluating}, the authors evaluate the applicability of OpenPose and MediaPipe Holistic for SLR using the ChaLearn dataset \cite{sincan2021chalearn}. They focus on interpretability through human-grounded evaluation \cite{doshi2017towards}, visually analyzing skeletal pose outputs. In \cite{raihan2024bengali}, a CNN is used alongside the KU-BdSL dataset \cite{jim2023ku}, with interpretability assessed using SHapley Additive eXplanation \cite{lundberg2017unified} to emphasize hand-related visual cues. In \cite{hu2023self}, a combination of 2D and 1D CNNs with a Bidirectional LSTM is employed for continuous SLR, using the PHOENIX14 \cite{camgoz2018neural}, PHOENIX14T \cite{koller2015continuous}, CSL-Daily \cite{zhou2021improving}, and CSL \cite{huang2018video} datasets, where interpretability is explored through Grad-CAM \cite{selvaraju2017grad} visualizations, showcasing the network's focus on crucial spatial features.

Interpretability in this field is largely underexplored, and the majority of existing research focuses primarily on SLR, rather than on the broader and more complex task of SLT. Example of this are \cite{moryossef2021evaluating,raihan2024bengali,hu2023self}, where CNNs and LSTMs SLR models are interpreted through human-grounded evaluation \cite{doshi2017towards}, SHAP \cite{lundberg2017unified}, or Grad-CAM \cite{selvaraju2017grad}. As for SLT, while some studies offer qualitative examples of interpretability, we have not identified any previous research that conducts a systematic analysis of interpretability in this domain. \cite{zhang2023heterogeneous} introduces the Heterogeneous Attention Transformer model, an end-to-end encoder–decoder transformer evaluated on the PHOENIX2014T dataset, with self-attention maps used for interpretability to demonstrate an improved attention distribution across two examples. In \cite{yin2023gloss}, the Gloss Attention SLT Network is presented, also based on a transformer architecture and evaluated on the PHOENIX14T, CSL-Daily, and SP-10 \cite{yin2022mlslt} datasets, with self-attention map visualizations shown only for a single example. Finally, in \cite{tang2022gloss}, the Gloss semantic-Enhanced Network with Online Back-Translation is proposed, integrating a gloss encoder, a pose decoder, and an online reverse gloss decoder for semantic alignment. Interpretability in the PHOENIX14T dataset is examined through cross-attention mechanism visualizations, but only two examples are provided to ensure consistency between gloss and pose sequences. Instead of providing simple qualitative examples, here we present a comprehensive interpretability analysis of a Transformer-based SLT model, shedding light on the mechanisms that allow the interaction between both sign and written languages.
